What are the random walk constructions of Brownian motion, arithmetic Brownian
motion and Brownian motion with time-dependent, but deterministic drift and
volatility and their multivariate counterparts? 
How does the number of operations depend on the number of points of
the time grid?

A first approach exploits the fact that the increments are normally distributed and independent.
On a time grid, process can thus be adjusted on the basis of these increments.

We fix times \(0=t_{0}<t_{1}<\cdot\cdot\cdot<t_{n}\) and independent random variables \(Z_{1},Z_{2},..., Z_{n}\sim\mathcal{N}(0,1).\)
1. Brownian motion
\(\mathcal{W}(t_{i+1})=W(t_{i})+\sqrt{t_{i+1}-t_{i}}\cdot Z_{i+1},\ i=0,1,...,n-1\)
2. Constant drift and volatility
\(X(t_{i+1})=X(t_{i})+\mu\cdot(t_{i+1}-t_{i})+\sigma\cdot\sqrt{t_{i+1}-t_{i}}\cdot Z_{i+1},\ i=0,1,...,n-1\)
3. Time-dependent drift and volatility
\(X(t_{i+1})=X(t_{i})+\int_{t_{i}}^{t_{i+1}}\mu(s)ds+\sqrt{\int_{t_{i}}^{t_{i+1}}\sigma^{2}(u)du}\cdot Z_{i+1},\ i=0,1,...,n-1\)

The random walk construction provides a simulation method that is exact on the discrete time grid
\(0=t_{0}<t_{1}<\cdot\cdot\cdot<t_{n}\)

For intermediate time points the process may be interpolated linearly. At these times the result is subject to
discretization error.
This applies, of course, also to all other methods that are exact on a time grid, but interpolate between the
specified time points.

The run-time complexity dependent on the resolution \( n \) of the grid is \( \mathcal{O}(n) \).

Fromn this construction one can also build the arithmetic Brownian motion by inserting it into the discretization
of its defining equation
\(dX_t = \mu dt + \sigma dW_t \).

The multivariate generalizations are given by
fixing times \(0=t_{0}<t_{1}<\cdot\cdot\cdot<t_{n}\) and independent random vectors \(Z_{1},Z_{2},...,Z_{n}\sim\mathcal{N}(0,I)\)
Then for \( i=0,1,...,n-1\) the schemes are
Brownian motion W:
\(W(t_{i+1})=W(t_{i})+\sqrt{t_{i+1}-t_{i}}\cdot Z_{i+1}\)

Constant drift and volatility \(dX(t)=\mu~dt+BdW(t)\):

\(X(t_{i+1})=X(t_{i}) + \mu \cdot (t_{i+1}-t_{i}) + \sqrt{t_{i+1}-t_{i}} \cdot B \cdot Z_{i+1}\)

Time-dependent drift and volatility \(dX(t)=\mu(t)dt+B(t)dW(t)\)

\(X(t_{i+1})=X(t_{i})+\int_{t_{i}}^{t_{i+1}}\mu(s)ds+B(t_{i},t_{i+1})\cdot Z_{i+1},\)
\(B(t_{i},t_{i+1})B(t_{i},t_{i+1})^{\top}=\int_{t_{i}}^{t_{i+1}}B(u)B(u)^{\top}du,\)

Derive the auto-covariance of Brownian motion. How does the Cholesky factor-
ization for finite-dimensional marginal distributions of Brownian motion look like?
When using this for simulation, how does the number of operations depend on
the number of points of the time grid, if the method is implemented directly (i.e.,
without additional simplifications)?

The Cholesky factorization provides an alternative simulation method that is based on the following observation:

\[ \text{Cov}(W(s), W(t)) = \text{Cov}(W(s), W(s) + (W(t) - W(s))) \]
\[ = \text{Cov}(W(s), W(s)) + \text{Cov} (W(s), W(t) - W(s)) = s \text{ with } s \leq t \]

Letting \(C \in \mathbb{R}^{n \times n}\) with \(C_{ij} = \min(t_i, t_j)\), we get \((W(t_1), W(t_2), ..., W(t_n)) \sim \mathcal{N}(0, C)\).

The Cholesky factorization \(C = AA^T\) is specified by
\[ A = 
\begin{pmatrix}
\sqrt{t_1} & 0 & \dots & 0 \\
\sqrt{t_1} & \sqrt{t_2 - t_1} & 0 & \vdots \\
\vdots & \vdots & \ddots & 0 \\
\sqrt{t_1} & \sqrt{t_2 - t_1} & \dots & \sqrt{t_n - t_{n-1}}
\end{pmatrix}
\] 

Brownian motion with drift \(\mu\) and volatility \(\sigma^2\) possesses on the time grid \(0=t_0<t_1<t_2<...<t_n\) a multivariate normal distribution
\(
N \left( \begin{pmatrix}
t_1 \ t_2 \ \vdots \ t_n
\end{pmatrix} \cdot \mu, \sigma^2 \cdot C \right)
\)

Using the Cholesky factorization, this can be simulated on the basis of a random vector Z with n independent standard normal components:
\(
(t_1 \ t_2 \ \dots \ t_n)^T \cdot \mu + \sigma \cdot A Z
\)


Describe the Brownian bridge simulation of Brownian motion. Where do you need
to argue with the Markov property? How does the number of operations depend
on the number of points of the time grid?

Benefit: later useful for variance reduction techniques and low-discrepancy methods
Connection to Brownian Bridge:

A Brownian motion conditioned on end points is a Brownian bridge:
1.generate \(W(t_n)\)
2. generate \(W(t_{\frac{n}{2}})\)
3. etc.

The specific implementation requires simulations from conditional distributions of three-dimensional random vectors.


We begin with considering a three-dimensional marginal distribution of Brownian motion, i.e.,
\( 0 \leq u < s < t \)

\(
\begin{pmatrix}
W(u) \\ W(s) \\ W(t)
\end{pmatrix} 
\sim \mathcal{N}
\left(
\begin{pmatrix}
0 \\ 0 \\ 0
\end{pmatrix},
\begin{pmatrix}
u & u & u \\
u & s & s \\
u & s & t
\end{pmatrix}
\right) 
\)
From the conditioning formula for normal random vectors we obtain:

\[\mathcal{L}(W(s) | W(u) = x, W(t) = y) = \mathcal{N} \left( \frac{(t-s)x + (s-u)y}{t-u}, \frac{(t-s)(s-u)}{t-u} \right) \]
This is the only formula that is needed for the Brownian bridge construction:
If the values
\[ \frac{W(s_1),}{= x_1} \frac{W(s_2),..., \frac{W(s_k)}{= x_k}} \]
are known for \(s_1 < s_2 < ... < s_k\), then for \(s \in (s_i, s_{i+1})\) conditioning on all values is the same as conditioning on \(W(s_i) = x_i, W(s_{i+1}) = x_{i+1}\) due to the Markov property of Brownian motion.

Sampling \(W(s)\) with \(s_i < s < s_{i+1}\), \(W(s_i) = x_i\), \(W(s_{i+1}) = x_{i+1}\):

\( W(s) = \frac{(s_{i+1} - s)x_i + (s - s_i)x_{i+1}}{s_{i+1} - s_i} + \sqrt{\frac{(s_{i+1} - s)(s - s_i)}{s_{i+1} - s_i}} \cdot Z, \quad Z \sim N(0,1) \)
* Sampling with drift
▸ sample rhs end point first
▸ afterwards the algorithm is identical, since a Brownian bridge constructed from Brownian motion with or without drift has the same law
* Sampling with volatility
▸ adjust the standard deviation of the innovations accordingly

The run-time complexity dependent on the resolution \( n \) of the grid is \( \mathcal{O}(n) \).

How could a principle component analysis be implemented for any multivariate
Gaussian process?

We consider the n-dimensional marginal distribution of a Brownian motion:
\( i, j \in \{1, 2, \dots, n\}   \)
\( (W(t_1) \ W(t_2) \ ... \ W(t_n))^T \sim \mathcal{N}(0, C) \text{ with } C_{ij} = \min(t_i, t_j), \ \)

* Solving 
\[ Cv_i = \lambda_i v_i, \ ||v_i|| = 1, \ i = 1, 2, ..., n \]
\[ \lambda_1 > \lambda_2 > ... > \lambda_n > 0 \]
\[ a_i = \sqrt{\lambda_i} v_i, \ i = 1, 2, ..., n \]
provides for suitable independent \( Z_i \sim \mathcal{N}(0,1), \ i = 1, 2, ..., n \) the representation
\[ (W(t_1) \ W(t_2) \ ... \ W(t_n))^T = Z_1 a_1 + Z_2 a_2 + ... + Z_n a_n \]

If \( t_{i+1} - t_i = \Delta t \ \forall i \), the solution is
\[ v_i(j) = \frac{2}{\sqrt{2n+1}} \cdot \sin \left( \frac{2i-1}{2n+1} \cdot j \pi \right), \ j = 1, 2, ..., n \]
\[ \lambda_i = \frac{\Delta t}{4} \sin^{-2} \left( \frac{2i-1}{2n+1} \cdot \frac{\pi}{2} \right) \]

The run-time complexity dependent on the resolution \( n \) of the grid is \( \mathcal{O}(n^{2}) \),
but the principle components construction identifies what is most important:
This knowledge can be exploited, e.g. to provide focus when variance reduction techniques are designed and
combined with the underlying approach.

State some general information about the geometric Brownian motion (gBM)
provide an exact recursive simulation scheme with deterministic drift and volatility
and the asymptotic behavior of the paths.

By the exponentiation of the arithmetic BM and applying Ito's lemma
the gBM satisfies the SDE
\( \frac{dS_t}{S_t} = (\mu + \frac{1}{2}\sigma^2)dt + \sigma dW_t \)
with the solution \( S_t = S(0)\exp((\mu \frac{1}{2}\sigma^2)t + \sigma W(t) \)

The one-dimensional marginal distributions of Geometric Brownian Motion are lognormal:
\(Y \sim LN(\mu, \sigma^2)\), if \(Y = \exp(\mu + \sigma Z)\), \(Z \sim N(0,1)\)
distribution function: 
([ P(Y \le y) = \Phi \left( \frac{\log(y) - \mu}{\sigma} \right) \]
 density: \( y \mapsto \frac{1}{y \sigma} \phi \left( \frac{\log(y) - \mu}{\sigma} \right) \)
 moments: 
\( \mathbb{E}(e^{aZ}) = e^{\frac{1}{2} a^2}, \ \mathbb{E}(Y) = e^{\mu + \frac{1}{2} \sigma^2}, \ \text{Var}(Y) = e^{2\mu + \sigma^2} (e^{\sigma^2} - 1) \)
median: \(\text{median}(Y) = e^\mu\)
In terms of the parameters of the Geometric Brownian Motion this translates as follows:
If \(S \sim GBM(\mu, \sigma^2)\), then \(\frac{S(t)}{S(0)} \sim LN((\mu - \frac{1}{2} \sigma^2)t, \sigma^2 t)\)
moments: \(\mathbb{E}(S(t)) = e^{\mu t} S(0), \ \text{Var}(S(t)) = e^{2 \mu t} S(0)^2 (e^{\sigma^2 t} - 1)\)
\( \mathbb{E}(S(t) | S(\tau), 0 \le t \le u) = \mathbb{E}(S(t) | S(u)) = e^{\mu (t - u)} S(u) \) due to the Markov property
* median: \(\text{median}(S(t)) = S(0) \cdot e^{(\mu - \frac{1}{2} \sigma^2)t} \)


The asymptotic behavior is given by
\( \frac{1}{t}W(t) \to 0 \) as \( t \to \infty \) a.s. by LLN so 
\( \frac{1}{t}\log(S_t) \to \mu - \frac{1}{2}\sigma^{2} \) a.s.
In particular if \( \frac{1}{2}\sigma^{2} > 0 \) then \( S_t \to \infty \) otherwise to \( 0 \).
Observe that 
\( \mu > 0 > \mu - \frac{1}{2}\simga^2 \implies S_t \to 0,\ \mathbb{E}[S_t] = e^{\mu t}S_0 \to \infty \)
a.s. as \( t \to \infty \).
Rare but very large values generate the exponential growth of the mean, although the paths converge to 0

Consider a time grid \(0 = t_0 < t_1 < ... < t_n\).
Random inputs to the simulation are independent \(Z_1, Z_2, ..., Z_n \sim \mathcal{N}(0, 1)\).
The following recursion is an exact simulation, i.e., without any discretization error:
for \( i = 1, 2, ..., n-1 \)
\( S(t_{i+1}) = S(t_i) \cdot \exp \left( (\mu - \frac{1}{2} \sigma^2) (t_{i+1} - t_i) + \sigma \sqrt{t_{i+1} - t_i} \cdot Z_{i+1} \right)\)
A generalized recursion is also available for time-dependent \(\mu\) and \(\sigma\):
\[ S(t_{i+1}) = S(t_i) \cdot \exp \left( \int_{t_i}^{t_{i+1}} (\mu(s) - \frac{1}{2} \sigma^2(s)) ds + \sqrt{\int_{t_i}^{t_{i+1}} \sigma^2(u) du} \cdot Z_{i+1} \right)\]


If dividends are introduced in the Black-Scholes-model, how do the accounting and
the risk-neutral measure have to be adjusted?

A continuous dividend yield \(\delta\) is associated with a dividend account:
\(\frac{dD(t)}{dt} = \delta S(t) + r D(t)\)
Need adjusted martingale condition, since now the process
\(\beta(t)^{-1} \cdot (S(t) + D(t)),\ t \geq 0 \)
is a martingale under \( Q \) the risk neutral measure
\(\Rightarrow \frac{dS(t)}{S(t)} = (r - \delta) dt + \sigma dW(t) \),
where \( \beta(t) = e^{rt} \) is a savings account

This is important to compute the paths of equity indices, which typically
include dividends in their growth as well as exchange rates, where 
the foreign interest rate can be seen as dividens \( r_f \)
and last commodities, where the costs of storing are negative dividends
and the benefits of holding are positive dividends, also called convenience yield.

Provide a couple of path-dependent options whose pricing typically requires simu-
lation.

Asian options
The payoff is the \( (\bar{S} - K)^{+} \) for a call and \( (K - \bar{S})^{-} \) for a put where \( \bar{S} \) is the average 
asset price during the considered time period.

Conventionally, this means an arithmetic average. In the continuous case, this is obtained by
\( \bar{S}(0,T) = \frac{1}{T} \int_{0}^{T} S(t) dt \)

For the case of discrete monitoring (with monitoring at the times \( 0=t_0, t_1, t_2, \dots, t_n=T \) and \( t_i=i\cdot \frac{T}{n} \))  we have the average given by
\( \bar{S}(0,T) = \frac{1}{n} \sum_{i=1}^{n} S(t_i) \)

Barrier options:
These exists in the following variations with barrier \( b \), strike \( K \) and expiration \( T \)
Up-and-out: spot price starts below the barrier level and has to move up for the option to be knocked out.
Down-and-out: spot price starts above the barrier level and has to move down for the option to become null and void.
Up-and-in: spot price starts below the barrier level and has to move up for the option to become activated.
Down-and-in: spot price starts above the barrier level and has to move down for the option to become activated.

Knock out means becoming worthless
so the payoff of an down-and-out could be described as 
\( \mathbb{1}_{\gamma(b) > T}(S(T) - K)^{+} \)
with \( \gamma(b) = \inf_{S_{t_i} < b} t_i \) (discr.)
or \( \gamma(b) = \inf_{S_{t_i} \leq b} t_i \) (cont.)

Lookback options
The option's strike price is floating (in the lecture there are also variations with a fixed strike \( K \))
and determined at maturity. 
The floating strike is the optimal value of the underlying asset's price during the option life. 
The payoff is the maximum difference between the market asset's price at maturity and the floating strike. 
For the call, the strike price is fixed at the asset's lowest price during the option's life, and, for the put, 
it is fixed at the asset's highest price.
So the payoff for a call is
\( \sup_{t \in [0, T]} S_T - S_t = S_T - S_{\text{min}}\)
and for a put
\( \sup_{t \in [0, T]} S_t - S_T  = S_{\text{max}} - S_T \)
Depending on cont./discr. the points admissable might not be the whole interval \( [0, T] \).

Explain that deterministic, time-dependent short-rates and volatilities still admit
an exact simulation of the stock. Why is this not true for a state-dependent local
volatility model?

Zero coupon bonds pay off \( 1 \) unit of currency at their maturity \( T \)
so their value in between is given by \( B(0, t) = \exp(-\int_{0}^{t} r(u)du \)
where \( r(u) \) is the risk-free rate, which could also be determined from
the bond prices via \( r(u) = -\frac{d}{dT}\log B(0, T)|_{T=u} \).
This requires the following adjustment to the price of the stock
\( S(t) = S(0)\exp(\int_{0}^t r(u)du - \frac{1}{2}\sigma^2 t + \sigma W(t) \)

We consider a time grid: \(0=t_{0}<t_{1}<\cdot\cdot\cdot<t_{n}\) and independent \(Z_{1},Z_{2},..., Z_{n}\sim\mathcal{N}(0,1)\)
Then
\(S(t_{i+1})=S(t_{i})\cdot exp \left( \int_{t_{i}}^{t_{i+1}}r(u)du-\frac{1}{2}\sigma^{2}(t_{i+1}-t_{i})+\sigma\sqrt{t_{i+1}-t_{i}}\cdot Z_{i+1} \right)\)
Alternatively, this can be rewritten in terms of zero-coupon bonds. This is due to the following identity:

\(B(t_{i},t_{i+1})=\frac{B(0,t_{i+1})}{B(0,t_{i})}=exp \left( -\int_{t_{i}}^{t_{i+1}}r(u)du \right)\)
The exact simulation can thus be rewritten in the following way:
\(S(t_{i+1})=S(t_{i})\cdot\frac{B(0,t_{i})}{B(0,t_{i+1})}\cdot exp \left( -\frac{1}{2}\sigma^{2}(t_{i+1}-t_{i})+\sigma\sqrt{t_{i+1}-t_{i}}\cdot Z_{i+1} \right)\)

If however the volatility is also dependent on the current price no exact simulation is available (only with time it should be, though not considered in the lecture)
and in inexact scheme with discretization errors is given by
\( S_{t_{i+1}} = S_{t_i}\exp((r - \frac{1}{2}\sigma^2(S_{t_i},t_i))(t_{i+1} t_i) + \sigma(S_{t_i}, t_i)\sqrt{t_{i+1} - t_i}Z_{i+1} \)


How can a multi-variate correlated geometric Brownian motion with fixed correla-
tion matrix be recursively simulated using a Cholesky decomposition?

Let  \(X_1, X_2, ..., X_d\) be standard Brownian motions that are correlated, i.e.,

\[\text{Cov}(X_i(t), X_j(t)) = \rho_{ij} t \quad \forall i,j,t\]
Setting \(\Sigma_{ij} = \sigma_i \sigma_j \rho_{ij}\), we introduce the notation
\[(\sigma_1 X_1 \ \sigma_2 X_2 \dots \ \sigma_d X_d)^T \sim BM(0, \Sigma)\]
We determine A such that \(AA^T = \Sigma\), e.g., as a Cholesky factor.
These correlated Brownian motions provide a basis for defining correlated Geometric Brownian motions:
\[\frac{dS_i(t)}{S_i(t)} = \mu_i dt + \sigma_i dX_i(t) = \mu_i dt + \sum_{j=1}^d A_{ij} dW_j(t),\ i = 1, 2, ..., d\]
for suitable uncorrelated Brownian motions \(W_1, W_2, ..., W_d\). One computes

\[\text{Cov}(S_i(t), S_j(t)) = S_i(0) S_j(0) e^{(\mu_i + \mu_j)t} (e^{\rho_{ij} \sigma_i \sigma_j t} - 1)\] 

Consider again a time grid \(0=t_{0}<t_{1}<\dots<t_{n}\) and independent d-dimensional standard Gaussian vectors
\[Z_{1},Z_{2},\dots,Z_{n}\sim\mathcal{N}(0,I_{d})\]

An exact simulation of the Geometric Brownian Motion in multiple dimensions is specified as follows:

\[S_{i}(t_{k+1})=S_{i}(t_{k})\cdot \exp \left( (\mu_{i}-\frac{1}{2}\sigma_{i}^{2})(t_{k+1}-t_{k})+\sqrt{(t_{k+1}-t_{k})\sum_{j=1}^{d}A_{ij}\cdot Z_{k+1,j}} \right),\]
for \( i \in \{ 1,2,\dots,d\} \)
and \( k \in \{ 0,1,2,\dots,n-1\} \)

The drift can be adjusted such that the simulation samples under \(P\), the statistical measure, or \(Q\), a risk-neutral measure, respectively.
If the savings account is the numeraire, then the drift under \(Q\) is the constant risk-free rate \(r\).

A multi-variate model is used for pricing of options that depend on multiple un-
derlyings. Provide examples of such options and their payoffs?

Spread Options
The spread between two assets at maturity is S1 (t) − S2 (T )
If the strike is K and we are considering a call on the spread, the corresponding payoff is
\( ([S_1(T) − S_2(t)] − K )^{+} \)

Basket Options
These options dependent on a portfolio of d underlying assets \( S_1, S_2, \dots, S_d \)
We suppose that the weights on the assets are \( c_1, c_2, \dots, c_d \)
A call with strike K on the basket has a payoff
\( ([\sum_{i=1}^{d} c_i S_i(T)] − K )^{+} \)

Outperformance options
An example of an outperformance option looks how a collection of assets outperforms a threshold \( K \)
This leads to a payoff
\( (\sup_{i} c_i S_i(T) -K)^{+}  \)
Generalized Barrier options
E.g.
A down-and-in put on \( S_1 \) that knocks in when \( S_2 \) drops below a barrier at \( b \):
\( \mathbb{1}_{\inf_{i \in \{1, \dots, n\} S_2(t_i) < b  }}(K - S_1(T))^{+} \)

Quantos
A quanto is a type of derivative in which the underlying is denominated in one currency, but the instrument
itself is settled in another currency at some rate.

Suppose e.g. that we are considering
the price S1 of a stock in a foreign currency
the exchange rate S2 expressed as the quantity of domestic currency required per unit of foreign currency
then the payoff is \( S_2(T)(S_1(T) - K)^{+} \).

Gaussian short rate models (generalizing the Ornstein-Uhlenbeck process) with
deterministic functions governing the drift and volatility admit a solution that can
conveniently be applied in special cases. State this.

The models are all special cases of an SDE of the following form
\( dr(t) = [g(t) + h(t)r(t)] dt + \sigma(t) dW(t) \)
with deterministic functions of time \(g\), \(h\), \(\sigma\)
The model admits the following solution:

\( r(t) = e^{H(t)} r(0) + \int_0^t e^{H(t)-H(s)} g(s) ds + \int_0^t e^{H(t)-H(s)} \sigma(s) dW(s), \ H(t) = \int_0^t h(s) ds \)
These processes are Gaussian, i.e., \(\forall t_1, t_2, ..., t_n\) the vector \((r(t_1), r(t_2), ..., r(t_n))\) is multivariate normal.

How do you derive in the Vasicek model a) an exact simulation and b) an Euler
scheme entailing discretization error?

Vasicek model
The model is based on an Ornstein-Uhlenbeck process
\( dr(t) = \alpha(b - r(t))dt + \sigma dW_t \)
with \( \alpha, b, \sigma > 0 \), \( W \) a \( Q \)-BM.
If the parameters are deterministic but time-dependent, the model is called the Hull-White model
We consider a version of the Vasicek model but with time-dependent b
\( r(t) = e^_{-\alpha(t-u)}r(u) + \alpha \int_{u}^{t} e^{-\alpha(t - s)} b(s) ds + \sigma \int_{u}^{t} e^{-\alpha(t-s}dW_s \)

Consider a time grid \(0=t_{0}<t_{1}<\dots<t_{n}\) and independent \(Z_{1},Z_{2},\dots,Z_{n}\sim\mathcal{N}(0,1)\).

We consider two simulation approaches. The first is based on the known conditional distribution; the second exploits naively the SDE.
1. Exact Simulation
\[r(t_{i+1})=e^{-\alpha(t_{i+1}-t_{i})}r(t_{i})+\mu(t_{i},t_{i+1})+\sigma_{r}(t_{i},t_{i+1})\cdot Z_{i+1}\]

2. Euler Scheme
\[r(t_{i+1})=r(t_{i})+\alpha(b(t_{i})-r(t_{i}))(t_{i+1}-t_{i})+\sigma\sqrt{t_{i+1}-t_{i}}\cdot Z_{i+1}\]

In contrast to the first method, the Euler scheme entails some discretization error and is not exact. 

Discuss calibration and simulation of the Ho-Lee model. Define forward rates and
explain their significance.

This model assumes a constant volatility \( \sigma \) but allows a flexible choice of the drift which is deterministic
function \( g \) of time:
\( dr(t) = g(t)dt + \sigma dW_t \)
Consider a time grid \(0=t_{0}<t_{1}<\dots<t_{n}\) and independent \(Z_{1},Z_{2},\dots,Z_{n}\sim\mathcal{N}(0,1)\).
 A generalized Ho-Lee model is characterized by the SDE
\[dX(t)=\mu(t)dt+\sigma(t)dW(t)\]
with deterministic functions \(\mu,\sigma\)

Exact Simulation
\[X(t_{i+1})=X(t_{i})+\int_{t_{i}}^{t_{i+1}}\mu(s)ds+\sqrt{\int_{t_{i}}^{t_{i+1}}\sigma^{2}(u)du}\cdot Z_{i+1}\]
Euler Scheme
\[X(t_{i+1})=X(t_{i})+\mu(t_{i})(t_{i+1}-t_{i})+\sigma(t_{i})\cdot\sqrt{t_{i+1}-t_{i}}\cdot Z_{i+1}\]

In contrast to the first method the Euler scheme entails some discretization error and is not exact. 
(Actually I think this is just a gBM)

Regarding the calibration, with some effort it can be seen
\[ g(t) = \frac{\partial}{\partial T} f(0, T)_{|T=t} + \sigma^2 t \]
where \( f(t, T) = - \frac{d}{dT}\log B(t, T) \) is the instantaneous forward rate observed at \( t \) for time \( T \),
satisfying \( B(0, T) = \exp(- \int_{0}^{T} f(0,t)dt) \):
Comparison with the bond price equation shows that

This allows to rewrite the SDE as
\[ dr(t) = g(t) dt + \sigma dW(t) = \left( \frac{\partial}{\partial T} f(0, T)_{|T=t} + \sigma^2 t \right) dt + \sigma dW(t) \]

This representation implies that, for simulating this model, g is not needed but only suitable forward rates.

Simulation
Consider a time grid \( 0 = t_0 < t_1 < ... < t_n \) and independent \( Z_1, Z_2, ..., Z_n \sim \mathcal{N}(0,1) \).
Exact simulation
\(r(t_{i+1}) &= r(t_i) + \int_{t_i}^{t_{i+1}} g(s) ds + \sigma \cdot \sqrt{t_{i+1} - t_i} \cdot Z_{i+1}\)
\(= r(t_i) + [f(0, t_{i+1}) - f(0, t_i)] + \frac{\sigma^2}{2} [t_{i+1}^2 - t_i^2] + \sigma \cdot \sqrt{t_{i+1} - t_i} \cdot Z_{i+1}\)

Recall a square-root diffusion and its simulation by Euler discretization. Show that
this simulation is not exact.

Any Gaussian process, like the Ornstein-Uhlenbeck process in the Vasicek model, takes negative values with positive probability.
This might be a disadvantage for interest rate and credit intensity models.
An interesting alternative goes back to William Feller (1951).
This process was introduced as a model for interest rates by John C. Cox, Jonathan E. Ingersoll and Stephen A. Ross in 1985.

Square-Root Diffusion (CIR-process)

\( dr(t) = \alpha \cdot (b-r(t))dt + \sigma \sqrt{r(t)} dW(t) \)

\( \alpha, b > 0, \ r(0) > 0, \ 2ab \geq \sigma^2 \)

Strong existence and uniqueness were shown by William Feller
This required some thought, since the square root is not Lipschitz-continuous at \( 0 \).

As simple methodology to simulate from a square-root diffusion is a Euler discretization.
We consider the points in time \(0 = t_0 < t_1 < ... < t_n\) and independent normals \(Z_1, Z_2, ..., Z_n \sim \mathcal{N}(0, 1)\).
The following Euler scheme is not exact:
\(r(t_{i+1}) = r(t_i) + \alpha (b - r(t_i)) [t_{i+1} - t_i] + \sigma \sqrt{r(t_i)^+} \sqrt{t_{i+1} - t_i} \cdot Z_{i+1}\)

A classical application is in the context of non-constant volatility models:

Assuming bond prices With the short rate
\( dr(t) = \alpha \cdot (b-r(t))dt + \sigma \sqrt{r(t)} dW(t) \)
they take the form
\( T \mapsto \mathbb{E} \left[ \exp \left( -\int_{0}^{T} r(u) du \right) \right] \)

This inspires the Heston model
This generalizes the Black-Scholes model by taking a square-root diffusion to specify the volatility of the asset:

\(\frac{dS(t)}{S(t)} = r dt + \sqrt{V(t)} dW_{1}(t)\)
\(dV(t) = \alpha \cdot (b-V(t)) dt + \sigma \sqrt{V(t)} dW_{2}(t)\)

The two-dimensional Brownian motion \((W_{1}, W_{2})\) may possess correlated components.


What is the density of a gamma distribution with given shape and scale parameter?
How do you obtain the special case of a \( \chi^{2} \)-distribution with \( \nu \) degrees of freedom?

Why is it sufficient to simulate gamma distributions with scale parameter equal to
1?


The gamma distribution is characterized by two parameters:
    * Shape parameter \(a\)
    * Scale parameter \(\beta\)
Its density is 
\(f(y) = f_{\alpha, \beta}(y) = \frac{1}{\Gamma(\alpha) \beta^\alpha} y^{\alpha-1} e^{-\frac{y}{\beta}}, y \ge 0\)
Example
A Chi-Square(\(\nu\)) corresponds to \(\beta = 2\) and \(\alpha = \frac{\nu}{2}\)
Simulation algorithms
Observe that we may focus on simulating \(X \sim \Gamma(\alpha, 1)\), since the scale parameter \(\beta\) captures that \(\beta X \sim \Gamma(\alpha, \beta)\).
We discuss two simulation methods, suggested by Cheng & Feast (1980) and Ahrens & Dieter (1974), that are appropriate for values \(a > 1\) and \(a \le 1\), respectively.

Elaborate on the different methods to simulate a gamma distribution

The method of Cheng and Feast builds on an acceptance-rejection method due
to Kinderman, Monahan. Explain the general framework of Kinderman,Monahan and prove that it works. 
What is the constraint on the shape parameter of the gamma distribution 
in the method of Cheng, Feast, and why is it assumed?

The approach of Cheng & Feast (1980) builds on a general methodology for random variable with values in \(\mathbb{R}_+\) and bounded density \(f\)
1. Generate \((U_1, U_2) \sim \text{unif}(A)\) with 
\(A = \{ (u_1, u_2) \in \mathbb{R}_+^2 : u_1 \le \sqrt{f \left( \frac{u_2}{u_1} \right) } \}\)
Then \(\frac{U_2}{U_1}\) has density \(f\).
2. In applications, \((U_1, U_2)\) are simulated by first generating bivariate random vectors that are uniformly distributed in a rectangle \(A\) (possibly except for some small mass) and then using the acceptance-rejection method.
This approach can be applied to a gamma distribution with \(\alpha > 1\).

For the gamma distribution with \(\alpha > 1\) one may use the rectangle \([0, 1] \times [0, \overline{u}_2]\) with

\(\overline{u}_1 = \left( \frac{\alpha - 1}{e} \right)^{\frac{\alpha - 1}{2}}, \quad \overline{u}_2 = \left( \frac{\alpha + 1}{e} \right)^{\frac{\alpha + 1}{2}},\)
which is a superset of
\(A = \left\{ (u_1, u_2) : 0 \leq u_1 \leq \sqrt{f \left( \frac{u_2}{u_1} \right)} \cdot e^{u_1} \right\}\)

 Applying the acceptance-rejection method allows to sample the gamma distribution \(\Gamma(\alpha, 1)\) from \(\frac{U_2}{U_1}\).

Repeat the following for the number of samples
\( \bar{a} \leftarrow a-1, \ b \leftarrow \left( a - \frac{1}{6a} \right) / \bar{a}, \ m \leftarrow \frac{2}{a}, \ d \leftarrow m+2 \)
Label A:
Generate \( U_1, U_2 \sim \text{unif}[0, 1] \)
\( V \leftarrow \frac{b U_2}{U_1} \) \\
if \( m U_1 - d + V + \frac{1}{V} \leq 0 \) goto label B
else if \( m \log U_1 - \log V + V - 1 \leq 0 \) goto label B
else goto Label A
Label B:
return \( Z \leftarrow \bar{a} V \)

The algorithm includes first an if-condition that does not require an evaluation of log; the evaluation of a more
complicated if-condition involving a log is only executed if necessary

For \( a \leq 1 \) one can use the acceptance-rejection method sampling first from

\( g(z) = p \cdot a z^{a-1}\) if \( z \in [0, 1] \)
otherwise if \( z > 1 \) \( g(z) = (1-p) \cdot e^{-z+1}\),
where \( p = \frac{e}{a + e} \)

The distribution with density \(g\) is a mixture:
with probability \(p\) generate a random variable on \([0, 1]\) with density \(z \mapsto a z^{a-1}\) using the inverse transform method
with probability \(1-p\) generate a random variable on \((1, \infty)\) with density \(z \mapsto e^{-z+1}\) using the inverse transform method
The density we are interested in is

\( f(z) = \frac{1}{\Gamma(a)} y^{a-1} e^{-y} \quad \text{with} \quad \frac{f(z)}{g(z)} \leq \frac{a + e}{ae \Gamma(a)},\ \forall z \)

The corresponding algorithmn is

Set \( b = \frac{a+e}{e} \)
Label A:
generate \( U_1, U_2 \sim \text{unif}[0, 1] \)
set \( Y = b U_1 \)
if \( Y \leq 1 \)
- set \( Z = Y^{1/a} \)
- if \( U_2 < e^{-Z} \) goto Label B
else \( Z = -\log(\frac{b - Y}{a} \)
if \( U_{2} \leq Z^{a-1} \) goto Label B
else goto Label A
Label B:
return \( Z \)

The correctness can be derived from
We prove that \((U_1, U_2) \sim \text{unif}(A)\) with \(A = \{ (u_1, u_2) \in \mathbb{R}^2_+ : u_1 \leq \sqrt{f \left( \frac{u_2}{u_1} \right) } \}\) implies that 
\(\frac{U_2}{U_1}\) has density \(f\).

\( \mathbb{P} \left( \frac{U_2}{U_1} \in B \right) = c \cdot \iint_{\mathbb{R}_+^2} 1_{\left\{ \frac{u_2}{u_1} \in B, (u_1, u_2) \in A \right\}} du_1 du_2 \)
\( = c \cdot \int_{\mathbb{R}_+} \int_{\mathbb{R}_+} 1_{\{ z \in B, (u_1, u_1 z) \in A \}} u_1 du_1 dz \)
\( = c \cdot \int_{\mathbb{R}_+} \int_{\mathbb{R}_+} u_1 \cdot 1_{\{ z \in B, u_1 \le \sqrt{f(z)} \}} du_1 dz \)
\( = \frac{c}{2} \int_B f(z) dz \)


Why is a recursive implementation of the inverse transform method particularly
efficient for the Poisson distribution?

A Poisson distribution \( \text{Poisson}(\theta) \) with parameter \(\theta > 0\) is given as follows:
\( N \sim \text{Poisson}(\theta) \)
\( \mathbb{P}(N = k) = e^{-\theta} \frac{\theta^k}{k!}, k = 0, 1, 2, ... \)
In the simulation we apply the inverse transform method.
The fact that \( \mathbb{P}(N = k) \) decays exponentially fast for \( k \to \infty \) suggests a recursive comparison.
An alternative approach based on exponential waiting times is slow:
\( X_1, X_2, ... \sim \text{exp}(\theta) \) with mean \(\frac{1}{\theta}\) independent
\( X_i = - \frac{\log(U_i)}{\theta} \) with \( U_1, U_2, ... \sim \text{unif}[0, 1] \) independent
\( N = \sup \{ n: X_1 + ... + X_n < 1 \} \) with \( \sup \emptyset = 0 \)

The algorithm then boils down to
\( p = e^{-\theta},\ F = p \)
\( N = 0 \)
generate \( U \sim \text{unif}[0, 1] \)
while \( U > F \)
\( N = N + 1 \)
\( p = \frac{p\theta}{N} \)
\( F = F + p \)
return \( N \)

